\documentclass[11pt,a4paper]{article}
\usepackage{auditstyle}
\lhead{Fresh Glimm--Jaffe proof audit}
\rhead{16 August 2026}
\usepackage{bm}

\newcommand{\T}{\mathbb T}
\newcommand{\W}{\mathcal W}
\newcommand{\eps}{\varepsilon}
\newcommand{\dd}{\,d}
\DeclareMathOperator{\gap}{gap}
\newtheorem{corollary}[theorem]{Corollary}

\title{Part II, corrected at the first falsifiable checkpoint\\[3pt]
\large The quadratic OAR model obstructs diagonal local state-norm convergence}
\author{Lamport-form proof and source audit}
\date{16 August 2026}

\begin{document}
\maketitle

\begin{resultbox}
\textbf{Decisive status.}  The Part II target in \texttt{strategy.md} is not a
well-defined statement as written, because an injective wavelet map embeds
\emph{observables}; it does not push a lattice state canonically to the continuum
algebra.  Under the natural MMST comparison, the intended diagonal local-norm
estimate is false already for the free massive scalar field.

More precisely, for every bounded interval $O$ containing a smaller interval, there
are finest-scale Weyl unitaries $U_N$ whose lattice support and continuum wavelet
support lie in $O$ and for which
\[
 \liminf_{N\to\infty}
 \left|\omega_N(U_N)-\omega\bigl(\beta_N(U_N)\bigr)\right|
 \geq e^{-\sqrt{2+\sqrt2}/4}-e^{-3\pi/16}
 =0.0752053795\ldots.                                            \tag{R}
\]
Consequently no estimate $C\eps_N^\theta$, $\theta>0$, can hold in local state
norm on the diagonal scale.  This does not contradict the MMST theorem: MMST take
the fine-scale limit while the tested observable is held at a \emph{fixed coarse
scale}.  Section \ref{sec:repair} gives the corrected Part II target.
\end{resultbox}

\tableofcontents

\section{Exact model and comparison map}\label{sec:model}

This section fixes every normalization.  The construction is the one-dimensional
specialization of Morinelli--Morsella--Stottmeister--Tanimoto (MMST), equations
(3.2)--(3.10), (4.1)--(4.5), and Proposition 4.1.

Fix $L>0$.  For $N\geq0$, choose integers $r_N=2^Nr_0$ and put
\[
 \eps_N=L/r_N,\qquad
 \Lambda_N=\{\eps_N n:-r_N\leq n\leq r_N-1\}\subset\T_L,
 \qquad \T_L=\R/(2L\mathbb Z).
\]
The dual momenta and their dimensionless representatives are
\[
 k_j=\frac{\pi j}{L},\qquad
 \theta_j=\eps_Nk_j=\frac{\pi j}{r_N},\qquad
 -r_N\leq j\leq r_N-1.                                         \tag{1.1}
\]
For a real sequence $p=(p_n)$, define
\[
 \widetilde p_j=\sum_{n=-r_N}^{r_N-1}p_ne^{-i\theta_jn}.
\]
Parseval's identity, with no suppressed constant, is
\[
 \sum_n|p_n|^2=\frac1{2r_N}\sum_j|\widetilde p_j|^2.             \tag{1.2}
\]

Let $\W_N$ be the Weyl algebra over $\C^{2r_N}$ with its standard symplectic
form.  In one space dimension the MMST decomposition (3.7) is
$\xi=\eps_Nq+ip$, $q,p$ real.  Let $m>0$ be the physical mass and put
$\delta_N=m\eps_N$.  The nearest-neighbour lattice dispersion is
\begin{align}
 \gamma_N(k_j)
 &=\left(m^2+2\eps_N^{-2}(1-\cos\theta_j)\right)^{1/2}\\
 &=\eps_N^{-1}a_{\delta_N}(\theta_j),
 &a_\delta(\theta)&=\left(\delta^2+4\sin^2(\theta/2)\right)^{1/2}. \tag{1.3}
\end{align}
The centered lattice vacuum is the quasifree state
\[
 \omega_N(W_N(\eps_Nq+ip))
 =\exp\left[-\frac14\left(
  \|\gamma_N^{-1/2}\widehat q\|_{N}^{2}
 +\|\gamma_N^{1/2}\widehat p\|_{N}^{2}\right)\right],          \tag{1.4}
\]
where $\widehat p_j=\eps_N^{1/2}\widetilde p_j$ and
$\|z\|_N^2=(2r_N)^{-1}\sum_j|z_j|^2$.

Let $\phi\in L^2(\R)$ be a real compactly supported scaling function such that
its integer translates are orthonormal.  We assume the exact decay used in MMST,
\[
 |\widehat\phi(\zeta)|\leq C(1+|\zeta|)^{-\rho},\qquad \rho>1,   \tag{1.5}
\]
which holds for the Daubechies choices admitted in MMST in one dimension.  Write
\[
\phi_{N,n}(x)=\eps_N^{-1/2}\phi(x/\eps_N-n).
\]
Periodize these functions with period $2L$.  The MMST symplectic map and the
induced injective morphism of Weyl $C^*$-algebras are
\begin{align}
 R_N(q,p)&=\eps_N^{1/2}\sum_n(q_n+ip_n)\phi_{N,n},               \tag{1.6}\\
 \beta_N(W_N(\eps_Nq+ip))&=W_{\rm ct}(R_N(q,p)).                 \tag{1.7}
\end{align}
Let $\omega$ be the continuum mass-$m$ vacuum on $\W(\T_L)$.

\begin{definition}[the only canonical finite-scale comparison]\label{def:comparison}
The lattice state $\omega_N$ and continuum state $\omega$ are compared on
$\W_N$ as
\[
 \omega_N\quad\hbox{versus}\quad \omega\circ\beta_N.           \tag{1.8}
\]
An injective $*$-homomorphism is contravariant on states: it sends $\omega$ to
$\omega\circ\beta_N$.  It supplies no canonical state on $\W(\T_L)$ from
$\omega_N$.
\end{definition}

\section{Exact covariance symbols}\label{sec:symbols}

\begin{lemma}[lattice momentum symbol]\label{lem:lattice-symbol}
For every real $p\in\R^{2r_N}$,
\[
 \omega_N(W_N(ip))
 =\exp\left[-\frac14 Q_N^{\rm lat}(p)\right],\qquad
 Q_N^{\rm lat}(p)=\frac1{2r_N}\sum_j
 a_{\delta_N}(\theta_j)|\widetilde p_j|^2.                      \tag{2.1}
\]
\end{lemma}

\subsection*{Proof in Lamport form}

\LS{1}{1}{Insert $q=0$ in (1.4).}

\LS{1}{2}{Use $\widehat p_j=\eps_N^{1/2}\widetilde p_j$ and
$\gamma_N(k_j)=\eps_N^{-1}a_{\delta_N}(\theta_j)$ to obtain}
\[
 \|\gamma_N^{1/2}\widehat p\|_N^2
 =\frac1{2r_N}\sum_j
  \eps_N^{-1}a_{\delta_N}(\theta_j)
  \eps_N|\widetilde p_j|^2,
\]
which is (2.1). \QedStep

\begin{lemma}[continuum pullback symbol]\label{lem:continuum-symbol}
For every real $p\in\R^{2r_N}$,
\[
 \omega(\beta_N(W_N(ip)))
 =\exp\left[-\frac14 Q_N^{\rm ct}(p)\right],                    \tag{2.2}
\]
where
\begin{align}
 Q_N^{\rm ct}(p)
 &=\frac1{2r_N}\sum_j b_{\phi,\delta_N}(\theta_j)
       |\widetilde p_j|^2,                                      \tag{2.3}\\
 b_{\phi,\delta}(\theta)
 &=\sum_{\ell\in\mathbb Z}
 \left(\delta^2+(\theta+2\pi\ell)^2\right)^{1/2}
 |\widehat\phi(\theta+2\pi\ell)|^2.                           \tag{2.4}
\end{align}
The series in (2.4) converges uniformly in $\theta\in[-\pi,\pi]$.
\end{lemma}

\subsection*{Proof in Lamport form}

\LS{1}{1}{Compute the Fourier coefficient of the embedded momentum smearing.}

For continuum momentum $k\in(\pi/L)\mathbb Z$, (1.6) and
$\widehat{\phi_{N,n}}(k)=\eps_N^{1/2}e^{-ik\eps_Nn}
\widehat\phi(\eps_Nk)$ give
\[
 \widehat{\Im R_N(0,p)}(k)
 =\eps_N\widehat\phi(\eps_Nk)
   \sum_np_ne^{-ik\eps_Nn}.                                    \tag{2.5}
\]

\LS{1}{2}{Split every continuum momentum uniquely as}
\[
 k=k_j+\frac{2\pi\ell}{\eps_N},\qquad
 -r_N\leq j\leq r_N-1,\quad \ell\in\mathbb Z.                 \tag{2.6}
\]
Then $e^{-ik\eps_Nn}=e^{-i\theta_jn}$ and
\[
 \eps_N(m^2+k^2)^{1/2}
 =\left(\delta_N^2+(\theta_j+2\pi\ell)^2\right)^{1/2}.         \tag{2.7}
\]

\LS{1}{3}{Insert (2.5)--(2.7) into the continuum vacuum quadratic form.}

The torus Fourier normalization is $(2L)^{-1}\sum_k$.  Since
$\eps_N/(2L)=1/(2r_N)$, grouping the sum by (2.6) gives exactly
(2.3)--(2.4), with no limiting argument.

\LS{1}{4}{Prove convergence of (2.4).}

By (1.5), uniformly for $|\theta|\leq\pi$ and $|\ell|\geq1$,
\[
 \left(\delta^2+(\theta+2\pi\ell)^2\right)^{1/2}
 |\widehat\phi(\theta+2\pi\ell)|^2
 \leq C_{\delta} (1+|\ell|)^{1-2\rho}.
\]
The majorant is summable because $\rho>1$.  The Weierstrass test proves uniform
convergence. \QedStep

\begin{lemma}[alias partition of unity]\label{lem:alias}
For every $\theta\in\R$,
\[
 \sum_{\ell\in\mathbb Z}
 |\widehat\phi(\theta+2\pi\ell)|^2=1.                            \tag{2.8}
\]
Consequently, with $d(\theta)=\dist(\theta,2\pi\mathbb Z)$,
\[
 b_{\phi,\delta}(\theta)
 \geq h_\delta(\theta):=\left(\delta^2+d(\theta)^2\right)^{1/2}. \tag{2.9}
\]
\end{lemma}

\subsection*{Proof in Lamport form}

\LS{1}{1}{Use orthonormality of the integer translates.}

For every $n\in\mathbb Z$, Plancherel gives
\begin{align*}
 \delta_{n0}
 &=\int_\R\phi(x)\phi(x-n)\dd x\\
 &=\frac1{2\pi}\int_{-\pi}^{\pi}e^{in\theta}
   \left(\sum_{\ell\in\mathbb Z}
   |\widehat\phi(\theta+2\pi\ell)|^2\right)\dd\theta.         \tag{2.10}
\end{align*}
The periodic function in parentheses therefore has Fourier coefficients
$\delta_{n0}$ and equals $1$ almost everywhere.  Assumption (1.5) makes its
series uniformly convergent, hence continuous.  Equality almost everywhere between
continuous functions implies (2.8) everywhere.

\LS{1}{2}{For each $\ell$,}
\[
 |\theta+2\pi\ell|\geq d(\theta),
\]
so every square-root factor in (2.4) is at least $h_\delta(\theta)$.  Multiply by
the nonnegative weights in (2.8) and sum.  This proves (2.9). \QedStep

\begin{remark}[strict ultraviolet mismatch]
For $0<|\theta|<\pi$ and $\delta\geq0$,
\[
 h_\delta(\theta)=\sqrt{\delta^2+\theta^2}
 >\sqrt{\delta^2+4\sin^2(\theta/2)}=a_\delta(\theta),            \tag{2.11}
\]
because $2\sin(|\theta|/2)<|\theta|$.  No property of a particular Daubechies
filter can remove this inequality.
\end{remark}

\section{A localized high-frequency sequence}\label{sec:carrier}

Fix open intervals $J\Subset O\Subset\T_L$ and choose real
$\eta\in C_c^\infty(J)$ with $\int_{\T_L}\eta(x)^2\dd x=1$.
Set
\[
 \theta_0=\frac{3\pi}{4},\qquad
 p_{N,n}=c_N\sqrt{2\eps_N}\,\eta(\eps_Nn)\cos(\theta_0n),        \tag{3.1}
\]
where $c_N>0$ is chosen so that $\sum_n p_{N,n}^2=1$.  For all
sufficiently large $N$, $2r_N$ is divisible by $8$, so the carrier is periodic.

\begin{lemma}[localized carrier measure]\label{lem:carrier}
For every continuous even $2\pi$-periodic function $f$,
\[
 \lim_{N\to\infty}\frac1{2r_N}\sum_j
 f(\theta_j)|\widetilde p_{N,j}|^2=f(\theta_0).                  \tag{3.2}
\]
Moreover $c_N\to1$.
\end{lemma}

\subsection*{Proof in Lamport form}

\LS{1}{1}{Prove $c_N\to1$.}

Before multiplication by $c_N^2$, the squared norm in (3.1) equals
\begin{align}
 2\eps_N\sum_n\eta(\eps_Nn)^2\cos^2(\theta_0n)
 &=\eps_N\sum_n\eta(\eps_Nn)^2\\
 &\quad+\eps_N\sum_n\eta(\eps_Nn)^2\cos(2\theta_0n).           \tag{3.3}
\end{align}
The first sum tends to $\int\eta^2=1$ by Riemann integration.  For the second,
partition $n$ into its four residue classes modulo $4$.  Since
$2\theta_0=3\pi/2$, each residue-class Riemann sum tends to
$\frac14\int\eta^2$, while
\[
 \sum_{a=0}^{3}e^{i(3\pi/2)a}=0.                                \tag{3.4}
\]
Thus the second sum tends to zero.  The unnormalized squared norm tends to one,
so $c_N\to1$.

\LS{1}{2}{Prove (3.2) first for $f(\theta)=e^{is\theta}$,
$s\in\mathbb Z$.}

By the discrete Fourier inversion formula, the left side of (3.2) is
\[
 \sum_np_{N,n}p_{N,n-s},                                        \tag{3.5}
\]
where no wrap-around term occurs for large $N$ because $\supp\eta\Subset(-L,L)$.
The product-to-sum identity turns (3.5) into
\begin{align}
 c_N^2\eps_N\sum_n&\eta(\eps_Nn)\eta(\eps_N(n-s))
 \left[\cos(s\theta_0)+\cos(2n\theta_0-s\theta_0)\right].      \tag{3.6}
\end{align}
The first term tends to $\cos(s\theta_0)$.  For the second, split $n=4q+a$,
$a=0,1,2,3$.  On each residue class the phase
$e^{i(2n\theta_0-s\theta_0)}$ equals
$e^{-is\theta_0}e^{i(3\pi/2)a}$, while the corresponding Riemann sum tends to
$\frac14\int\eta(x)^2dx$ because
$\eta(x)\eta(x-\eps_Ns)\to\eta(x)^2$ uniformly.  Summing the four limits gives
zero by (3.4).
Therefore (3.5) tends to
\[
 \cos(s\theta_0)=\frac12(e^{is\theta_0}+e^{-is\theta_0}).       \tag{3.7}
\]

\LS{1}{3}{Extend to continuous $f$.}

Equation (3.7) proves weak convergence of the probability measures
\[
 \nu_N=\frac1{2r_N}\sum_j|\widetilde p_{N,j}|^2\delta_{\theta_j}
 \quad\Longrightarrow\quad
 \frac12(\delta_{\theta_0}+\delta_{-\theta_0})                 \tag{3.8}
\]
first on trigonometric polynomials.  Trigonometric polynomials are uniformly dense
in the continuous periodic functions, and both sides of (3.8) have mass one.
Hence (3.8) holds for every continuous periodic $f$.  If $f$ is even, its integral
against the limiting measure is $f(\theta_0)$, proving (3.2). \QedStep

\section{The diagonal local-norm no-go theorem}\label{sec:nogo}

Let $X_N=\{\eps_Nn:p_{N,n}\neq0\}$ and let $\W_N(X_N)$ be the corresponding
lattice Weyl subalgebra.

\begin{theorem}[free quadratic obstruction]\label{thm:nogo}
For the MMST diagonal lattice vacua and continuum mass-$m$ vacuum,
\begin{align}
 \lim_{N\to\infty}Q_N^{\rm lat}(p_N)
 &=2\sin(3\pi/8)=\sqrt{2+\sqrt2},                               \tag{4.1}\\
 \liminf_{N\to\infty}Q_N^{\rm ct}(p_N)
 &\geq 3\pi/4.                                                   \tag{4.2}
\end{align}
Consequently
\begin{align}
 \liminf_{N\to\infty}
 \left\|
 \omega_N|_{\W_N(X_N)}-(\omega\circ\beta_N)|_{\W_N(X_N)}
 \right\|
 &\geq c_*,                                                       \tag{4.3}\\
 c_*&=e^{-\sqrt{2+\sqrt2}/4}-e^{-3\pi/16}>0.                    \tag{4.4}
\end{align}
Numerically, $c_*=0.0752053795\ldots$.
\end{theorem}

\subsection*{Proof in Lamport form}

\LS{1}{1}{Apply Lemma \ref{lem:carrier} to the lattice symbol.}

The functions $a_\delta$ converge uniformly to
$a_0(\theta)=2|\sin(\theta/2)|$, because
\[
 0\leq a_\delta(\theta)-a_0(\theta)\leq\delta.                  \tag{4.5}
\]
Equations (2.1), (3.2), and (4.5) yield
\[
 Q_N^{\rm lat}(p_N)\longrightarrow a_0(\theta_0)
 =2\sin(3\pi/8)=\sqrt{2+\sqrt2}.                                \tag{4.6}
\]

\LS{1}{2}{Apply the alias lower bound to the continuum symbol.}

By (2.9),
\[
 Q_N^{\rm ct}(p_N)
 \geq\frac1{2r_N}\sum_jh_{\delta_N}(\theta_j)
       |\widetilde p_{N,j}|^2.                                  \tag{4.7}
\]
The periodic functions $h_\delta$ converge uniformly to
$h_0(\theta)=d(\theta)$, because
$0\leq h_\delta-h_0\leq\delta$.  Lemma \ref{lem:carrier} and
$0<\theta_0<\pi$ therefore give
\[
 \liminf_NQ_N^{\rm ct}(p_N)\geq h_0(\theta_0)=3\pi/4.           \tag{4.8}
\]

\LS{1}{3}{Use the Weyl unitary itself as a norm-one test observable.}

By (2.1)--(2.2),
\[
 \left|\omega_N(W_N(ip_N))-\omega(\beta_N(W_N(ip_N)))\right|
 =\left|e^{-Q_N^{\rm lat}(p_N)/4}-e^{-Q_N^{\rm ct}(p_N)/4}\right|.
                                                                    \tag{4.9}
\]
The right side has lower limit at least (4.4), by (4.1)--(4.2) and the strict
inequality $\sqrt{2+\sqrt2}<3\pi/4$.  Since $W_N(ip_N)$ is a unitary in
$\W_N(X_N)$, the dual norm dominates (4.9).  This proves (4.3). \QedStep

\begin{corollary}[no continuum extension repairs the norm]\label{cor:extension}
Assume $\supp\phi\subset[0,R_\phi]$.  For all sufficiently large $N$,
\[
 \beta_N(\W_N(X_N))\subset\W(O).                                \tag{4.10}
\]
If $\widetilde\omega_N$ is any state on the continuum Weyl algebra satisfying
$\widetilde\omega_N\circ\beta_N=\omega_N$, then
\[
 \liminf_{N\to\infty}
 \| (\widetilde\omega_N-\omega)|_{\W(O)}\|\geq c_*.            \tag{4.11}
\]
\end{corollary}

\begin{proof}
Choose $a>0$ with $\dist(\supp\eta,\T_L\setminus O)>a$.  Every nonzero
coefficient in (3.1) is based at a point of $\supp\eta$, and
\[
 \supp\phi_{N,n}\subset\eps_Nn+[0,R_\phi\eps_N].               \tag{4.12}
\]
If $R_\phi\eps_N<a$, (4.10) follows.  On the subalgebra in (4.10), the two
functionals in (4.11) restrict to the two functionals in (4.3).  Restriction cannot
increase a functional norm, so (4.3) implies (4.11).
\end{proof}

\begin{corollary}[the proposed Part II vertical term is false]\label{cor:false-target}
There are no constants $C<\infty$ and $\theta>0$ for which
\[
 \left\|
 \omega_N|_{\W_N(X_N)}-(\omega\circ\beta_N)|_{\W_N(X_N)}
 \right\|\leq C\eps_N^\theta                                  \tag{4.13}
\]
holds for all sufficiently large $N$.
\end{corollary}

\begin{proof}
The right side of (4.13) tends to zero; the lower limit of the left side is at
least $c_*>0$ by Theorem \ref{thm:nogo}.
\end{proof}

\begin{remark}[spatial cutoffs do not evade the obstruction]
In the quadratic model, let $g_R\uparrow1$ and let the one-particle squared
energies be
\[
 A_{N,R}=-\Delta_N+m_0^2+\lambda g_R,\qquad
 B_N=-\Delta_N+m_0^2+\lambda.
\]
At fixed $N$, $A_{N,R}\to B_N$ strongly and
$m_0^2\leq A_{N,R},B_N\leq 4\eps_N^{-2}+m_0^2+\lambda$.  Continuous functional
calculus on this common compact spectral interval gives
$A_{N,R}^{\pm1/2}\to B_N^{\pm1/2}$ strongly.  Hence the expectation of every
finitely supported Weyl operator converges to the full mass-$m_1$ lattice vacuum,
$m_1^2=m_0^2+\lambda$.  Therefore a putative estimate
$\Psi(d_R)+C\eps_N^\theta$, with $\Psi(d_R)\to0$, would imply (4.13) after first
taking $R\to\infty$ at fixed $N$.  Theorem \ref{thm:nogo} rules this out.
\end{remark}

\section{The free facts that do survive (II.2)}\label{sec:freefacts}

\begin{proposition}[exact gap and group-velocity bounds]\label{prop:gap-velocity}
For the MMST free lattice Hamiltonian, after subtraction of the vacuum energy,
\[
 \gap H_N=m                                                       \tag{5.1}
\]
for every $N$.  Its one-particle group velocity satisfies
\[
 \sup_N\sup_{k\in[-\pi/\eps_N,\pi/\eps_N]}
 |\partial_k\gamma_N(k)|\leq1.                                  \tag{5.2}
\]
For the quadratic spatial-cutoff model
$A_N(g)=-\Delta_N+m_0^2+\lambda g$, $0\leq g\leq1$, the Fock gap is at least
$m_0$; on the infinite lattice with compactly supported $g$, it equals $m_0$.
\end{proposition}

\subsection*{Proof in Lamport form}

\LS{1}{1}{The one-particle energy is multiplication by $\gamma_N(k)$ and
$\gamma_N(k)\geq\gamma_N(0)=m$.}

The normal-ordered many-body Hamiltonian is
\[
 H_N=d\Gamma(\gamma_N).
\]
Its nonzero spectrum begins at $\inf\operatorname{spec}\gamma_N=m$, proving
(5.1).

\LS{1}{2}{Differentiate (1.3).}
\[
 \partial_k\gamma_N(k)
 =\frac{\sin(\eps_Nk)}
 {\sqrt{\delta_N^2+4\sin^2(\eps_Nk/2)}}.                         \tag{5.3}
\]
Since $|\sin\theta|=2|\sin(\theta/2)\cos(\theta/2)|$, the absolute value
of (5.3) is at most $|\cos(\theta/2)|\leq1$.  This proves (5.2).

\LS{1}{3}{For a spatial cutoff, $A_N(g)\geq m_0^2$, so
$A_N(g)^{1/2}\geq m_0$.}

If $g$ is compactly supported on the infinite lattice, multiplication by $g$ is a
compact perturbation of $-\Delta_N+m_0^2$.  Weyl's theorem gives essential
spectrum beginning at $m_0^2$.  The perturbation is nonnegative, so there is no
spectrum below $m_0^2$.  Thus the bottom is exactly $m_0^2$, and the Fock gap is
$m_0$. \QedStep

\begin{proposition}[exact Wick-constant asymptotic]\label{prop:wick}
The infinite-lattice equal-site field covariance is
\[
 c_N=\frac1{4\pi}\int_{-\pi}^{\pi}
 \frac{\dd\theta}{\sqrt{\delta_N^2+4\sin^2(\theta/2)}}.         \tag{5.4}
\]
As $N\to\infty$,
\[
 c_N=\frac1{2\pi}\log\frac{8}{m\eps_N}+o(1).                   \tag{5.5}
\]
In particular the coefficient of $\log(1/\eps_N)$ is exactly $1/(2\pi)$.
\end{proposition}

\subsection*{Proof in Lamport form}

\LS{1}{1}{Evenness and the substitution $\theta=2u$ give}
\[
 c_N=\frac1\pi\int_0^{\pi/2}
 \frac{\dd u}{\sqrt{\delta_N^2+4\sin^2u}}.                      \tag{5.6}
\]

\LS{1}{2}{Add and subtract the model singularity $2u$.}
\begin{align}
 c_N&=I(\delta_N)+D(\delta_N),                                  \tag{5.7}\\
 I(\delta)&=\frac1\pi\int_0^{\pi/2}
 \frac{\dd u}{\sqrt{\delta^2+4u^2}}
 =\frac1{2\pi}\operatorname{arsinh}\frac\pi\delta,           \tag{5.8}\\
 D(\delta)&=\frac1\pi\int_0^{\pi/2}\left[
 \frac1{\sqrt{\delta^2+4\sin^2u}}
 -\frac1{\sqrt{\delta^2+4u^2}}\right]\dd u.                   \tag{5.9}
\end{align}

\LS{1}{3}{Justify dominated convergence in (5.9).}

For $0<a\leq b$,
\[
 0\leq\frac1{\sqrt{\delta^2+a^2}}
       -\frac1{\sqrt{\delta^2+b^2}}
 \leq\frac1a-\frac1b,                                         \tag{5.10}
\]
because the absolute derivative of $(\delta^2+x^2)^{-1/2}$ is at most $x^{-2}$.
Apply (5.10) with $a=2\sin u$, $b=2u$.  The majorant
\[
 \frac1{2\sin u}-\frac1{2u}                                    \tag{5.11}
\]
is integrable at $u=0$, since it equals $u/12+O(u^3)$.  Therefore
\begin{align}
 \lim_{\delta\downarrow0}D(\delta)
 &=\frac1{2\pi}\int_0^{\pi/2}(\csc u-u^{-1})\dd u\\
 &=\frac1{2\pi}\log\frac4\pi.                                \tag{5.12}
\end{align}
The last equality follows by integrating from $a$ to $\pi/2$, using
$\int\csc u\dd u=\log\tan(u/2)$, and then taking $a\downarrow0$.

\LS{1}{4}{Since}
\[
 \operatorname{arsinh}(\pi/\delta)=\log(2\pi/\delta)+o(1),      \tag{5.13}
\]
adding (5.8) and (5.12) yields
\[
 c_N=\frac1{2\pi}\left[
 \log\frac{2\pi}{\delta_N}+\log\frac4\pi\right]+o(1)
 =\frac1{2\pi}\log\frac8{m\eps_N}+o(1).
\]
This proves (5.5). \QedStep

\begin{remark}[what (5.2) does and does not prove]
Equation (5.2) is a group-velocity estimate.  It is not, by itself, an interacting
Lieb--Robinson theorem for unbounded quartic oscillators.  MMST obtain a free harmonic
Lieb--Robinson estimate separately in their Lemma 4.14 and Proposition 4.15.  Part II
must import or prove an interacting bound with the exact $\eps_N$ scaling; relabeling
(5.2) as such a theorem would be invalid.
\end{remark}

\section{What the cited sources establish}\label{sec:sources}

\begin{proposition}[precise MMST scope]\label{prop:mmst-scope}
The following statements, and no stronger diagonal state-norm claim, are supported by
the cited MMST paper.
\begin{enumerate}[label=\textup{(\roman*)}]
\item The maps $\alpha_N^{N'}$ form an inductive system of observable algebras, while
states at finer scales are pulled back to a fixed coarse algebra.
\item Proposition 4.1 embeds the inductive-limit Weyl algebra into the continuum Weyl
algebra by $\beta_L$.
\item Lemmas 4.2--4.3 and Proposition 4.5 prove, for every \emph{fixed} coarse scale,
weak$^*$ convergence of the pulled-back free lattice ground states and identify the
projective limit with the continuum vacuum on the embedded algebra.
\item Theorem 4.11 identifies the resulting time-zero net with the massive free-field
net.
\item Section 6.1 explicitly treats interacting lattice ground-state scaling limits and
wavelet interacting dynamics as future generalizations.
\end{enumerate}
\end{proposition}

\begin{proof}
Items (i)--(iv) are the literal content of MMST equations (2.8)--(2.11),
Propositions 4.1 and 4.5, and Theorem 4.11.  Item (v) is stated in the first two
paragraphs of MMST Section 6.1.  Page images of all five locations appear in the
separate Part II citation dossier.
\end{proof}

\begin{proposition}[Battle--Federbush does not supply II.4]\label{prop:battle-scope}
The 1987 Battle--Federbush note does not state CE1$_N$/CE2$_N$ for the MMST
Daubechies OAR lattice, uniformly in $N$, nor for the mixed interpolation proposed in
II.6.
\end{proposition}

\begin{proof}
The note uses Meyer-type or exponentially localized orthonormal modes $\psi_k$ and
expands the Euclidean field with
\[
 u_k=(-\Delta+M^2)^{-1/2}\psi_k.                                \tag{6.1}
\]
It explicitly singles out the resulting diagonality of the free action in the phase-cell
variables as an attractive feature.  The MMST finite-scale Hamiltonian instead has the
nearest-neighbour symbol (1.3), and its Daubechies functions define symplectic scaling
maps between time-zero Weyl algebras.  The note contains no theorem with the MMST
Hamiltonian, the Wick constant (5.4), a sharp-time GJS observable norm, an $N$-uniform
physical gap, or two mixed vertex species.  These missing hypotheses are exactly the
conclusion demanded by II.4, so the note cannot be cited as that conclusion.

\end{proof}

\section{Why II.6 is not an operator interpolation}\label{sec:interpolation}

\begin{proposition}[ill-defined mixed Hamiltonian path]\label{prop:path}
The expression
\[
 H_s=(1-s)H_N(g)+sH(g)                                          \tag{7.1}
\]
in the original II.6 program is undefined with the currently specified data.
\end{proposition}

\begin{proof}
$H_N(g)$ is a self-adjoint operator in the lattice ground-state Hilbert space
$\Hh_{N,g}$, whereas $H(g)$ is a self-adjoint operator in the continuum Hilbert space
$\Hh_g$.  Addition of unbounded operators is defined only after both are realized as
operators or closed forms on one Hilbert space with a specified common form domain.
The observable embedding $\beta_N:\W_N\to\W$ supplies neither:
\begin{enumerate}[label=\textup{(\alph*)}]
\item it maps $C^*$-algebra elements, not unbounded Hamiltonians;
\item it does not push $\omega_N$ to a canonical continuum state;
\item its symplectic one-particle map need not intertwine the finite-scale and continuum
complex structures, so no stated second-quantized unitary identifies the two vacuum
representations;
\item no common closed quadratic form for the two Wick-ordered interactions is given.
\end{enumerate}
Therefore (7.1), its derivative $\partial_sH_s$, its spectral gap, and the vectors
$D_j\Omega_s$ have no definitions.  A mixed-vertex Euclidean measure is not a repair of
this operator defect unless an Osterwalder--Schrader construction identifies the entire
path on a common Hilbert space and proves common-domain control.
\end{proof}

\section{Corrected Part II target and dependency order}\label{sec:repair}

The no-go theorem forces a change of topology.  It does not merely change the value of
the exponent $\theta$.

\begin{definition}[fixed-coarse comparison]\label{def:fixed-coarse}
For $M<N$, a lattice state $\omega_{N,g}$ defines a state on the fixed coarse algebra
$\W_M$ by
\[
 \omega_{N,g}^{[M]}:=\omega_{N,g}\circ\alpha_M^N.               \tag{8.1}
\]
The continuum comparison state is
\[
 \omega_\infty^{[M]}:=\omega_\infty\circ\beta_M.                \tag{8.2}
\]
Both sides of (8.1)--(8.2) are states on exactly the same $C^*$-algebra.
\end{definition}

\begin{definition}[resolution-separated diagonal]\label{def:separated}
Let $M(N)\leq N$ satisfy
\[
 M(N)\to\infty,\qquad N-M(N)\to\infty.                          \tag{8.3}
\]
The first condition makes the tested observables dense in the continuum algebra; the
second prevents this comparison class from containing the particular finest-scale
$M=N$ carrier used in Theorem \ref{thm:nogo}.  These two conditions make the topology
admissible; they do not by themselves prove convergence.
\end{definition}

\begin{definition}[corrected interacting Part II target]\label{def:corrected-target}
For each fixed $M$, bounded lattice region $X\subset\Lambda_M$, and
$A\in\W_M(X)$,
\[
 \left|\omega_{N,g}^{[M]}(A)-\omega_\infty^{[M]}(A)\right|
 \leq \|A\|\,\Psi_\gamma(d)+E_{M,X,A}(N),                       \tag{8.4}
\]
where $E_{M,X,A}(N)\to0$ as $N\to\infty$, uniformly in the spatial cutoff
once $g=1$ on the required neighbourhood.  A stronger uniform estimate may be sought
along (8.3), but no supremum over all $M\leq N$ can tend to zero.
\end{definition}

\begin{proposition}[what is proved about the corrected target]\label{prop:target-status}
The left side of (8.4) is well-defined.  In the free full-space case MMST prove that it
tends to zero for each fixed $M$ and $A$.  No estimate (8.4) can hold uniformly over a
class containing $M=N$ with $E_{M,X,A}(N)\to0$ uniformly over its norm-one local
observables.
\end{proposition}

\begin{proof}
Definitions \ref{def:comparison} and \ref{def:fixed-coarse} put both functionals on
$\W_M$, proving the first assertion.  MMST Proposition 4.5 and Lemma 4.3 give the
second assertion.  For the third, take $M=N$, $X=X_N$, and the norm-one observable
$A=W_N(ip_N)$ of Theorem \ref{thm:nogo}.  The left side then has lower limit at least
$c_*>0$ by (4.3), contradicting any uniform error tending to zero.
\end{proof}

The corrected dependency order is therefore:
\begin{enumerate}[label=\textbf{II.\arabic*.},start=0]
\item \textbf{Choose one lattice model.}  Decide between the MMST nearest-neighbour
Hamiltonian, a Daubechies Galerkin Hamiltonian, or another action.  These choices have
different symbols, interactions, counterterms, and locality estimates.
\item \textbf{Checkpoint completed.}  The diagonal local-norm target fails by Theorem
\ref{thm:nogo}.  Fix either the fixed-coarse topology (8.4) or a precise
resolution-separated norm before proceeding.
\item \textbf{Free estimates partially completed.}  Equations (5.1)--(5.5) give the
MMST gap, group velocity, and Wick asymptotic.  A separate interacting
Lieb--Robinson theorem is still required.
\item \textbf{Source audit completed negatively.}  Battle--Federbush is mechanism
precedent, not CE1$_N$/CE2$_N$ for the chosen OAR action.
\item \textbf{Uniform interacting expansion remains open.}  State it directly in the
fixed-coarse or resolution-separated topology, including counterterms and source norms.
\item \textbf{LPPL with lattice tails remains open.}  It must be formulated on one
finite-scale algebra and must include quantitative commutator errors for affiliated
unbounded Wick polynomials.
\item \textbf{Replace the invalid vertical interpolation.}  Use a common Euclidean
interpolation followed by a proved OS identification, or compare projective states
directly.  Do not add $H_N$ and $H$ across different Hilbert spaces.
\end{enumerate}

\section{Claim ledger}

\begin{longtable}{P{0.27\textwidth}P{0.20\textwidth}P{0.41\textwidth}}
\toprule
Claim & Status & Exact reason\\
\midrule
\endhead
MMST supplies the free benchmark & \Pass & Fixed-coarse weak$^*$ projective limit;
Proposition 4.5 and Theorem 4.11.\\
The displayed Part II norm is defined & \Fail & An observable embedding has no
canonical push-forward on states.\\
Diagonal local state norm is $O(\eps_N^\theta)$ & \Fail & Explicit local Weyl lower
bound (4.3)--(4.4).\\
$N$-uniform free gap & \Pass & Exact identity $\gap H_N=m$.\\
$N$-uniform free group velocity & \Pass & Exact derivative bound (5.3).\\
Wick constant coefficient & \Pass & Explicit asymptotic (5.5).\\
Group velocity alone gives interacting LR & \Fail & It controls only the free
dispersion; an anharmonic commutator theorem is separate.\\
Battle supplies CE1$_N$/CE2$_N$ & \Fail & Different variables/action and no theorem
in the required norms or mixed family.\\
$H_s=(1-s)H_N+sH$ is defined & \Fail & Operators live on different Hilbert spaces;
no common form realization is specified.\\
Fixed-coarse interacting OAR estimate & \Conditional & Meaningful target; free case
known pointwise, interacting uniform expansion still absent.\\
Full Part II OAR closure as originally written & \Fail & False topology plus undefined
interpolation; it cannot be repaired by estimating a different exponent.\\
\bottomrule
\end{longtable}

\section{Reproducibility}

The script \texttt{m2\_ii\_checkpoint.py} evaluates the finite-torus quadratic forms
in (2.1) and the universal lower form from (2.9), the group-velocity bound (5.3), and
the Wick integral (5.4).  It asserts convergence to the exact constants in
(4.1)--(4.4) and (5.5).  The numerical output is diagnostic only; every limiting
statement used in the proof was derived analytically above.

\begin{thebibliography}{9}

\bibitem{MMST}
V. Morinelli, G. Morsella, A. Stottmeister, and Y. Tanimoto,
\emph{Scaling limits of lattice quantum fields by wavelets},
Commun. Math. Phys. \textbf{387} (2021), 1299--1367,
\href{https://doi.org/10.1007/s00220-021-04152-5}{doi:10.1007/s00220-021-04152-5}.

\bibitem{BF87}
G. Battle and P. Federbush,
\emph{Ondelettes and phase cell cluster expansions, a vindication},
Commun. Math. Phys. \textbf{109} (1987), 417--419,
\href{https://doi.org/10.1007/BF01206144}{doi:10.1007/BF01206144}.

\end{thebibliography}

\end{document}
